\newpage
\section{CUDA Parallelization}
\label{sec:parallelization}

\subsection{Sequential section}
I initially decided to run the tests on my personal PC rather than on the university workstation or on Google Colab, in order to avoid limitations related to temporary sessions and restricted computational resources. This allowed me to maintain full control over the compilation process and the execution of the tests during the initial experiments with CUDA.\\

\noindent
The simulation parameters used are:

\begin{itemize}
    \item $nstep = 200$
    \item $ni = 15000$
    \item $nj = 15000$
    \item $size = ni \times nj = 900\,000\,000$
\end{itemize}

\noindent
I compiled the sequential code located in \texttt{src/cuda.c} using the following options:

\begin{verbatim}
    gcc -O3 src/cuda.c -o build/cuda
\end{verbatim}

\noindent
The best execution time obtained (37 s) is used as the sequential baseline for computing the speedup of the GPU implementation.

\subsection{Code}
The parallelization was implemented using CUDA to execute the time-step computations on the GPU. The main kernel is \texttt{step\_kernel\_mod}, which replicates the behavior of the sequential function \texttt{step\_kernel\_ref}. You can find the code in \texttt{src/cuda.cu}.

\paragraph{Kernel structure:}
\begin{itemize}
    \item Each thread processes one point of the 2D grid.
    \item The global indices $i$ and $j$ are computed from \texttt{threadIdx} and \texttt{blockIdx}.
\end{itemize}

\paragraph{Memory management:}
\begin{itemize}
    \item Allocation of host and device memory.
    \item Transfer of the initial data from host to device.
    \item Transfer of the results from device to host at the end of the simulation.
\end{itemize}

\paragraph{Error checking:}
\begin{itemize}
    \item After each kernel launch, \texttt{cudaGetLastError()} is invoked.
    \item The maximum error relative to the CPU reference is computed to check correctness.
\end{itemize}

\begin{lstlisting}
__global__ void step_kernel_mod(int ni, int nj, float fact, float* temp_in, float* temp_out)
{
  int i00, im10, ip10, i0m1, i0p1;
  float d2tdx2, d2tdy2;

  int i = threadIdx.x + blockIdx.x * blockDim.x;
  int j = threadIdx.y + blockIdx.y * blockDim.y;

  if(i > 0 && i < ni-1 && j > 0 && j < nj-1){
    // find indices into linear memory
    // for central point and neighbours
    i00 = I2D(ni, i, j);
    im10 = I2D(ni, i-1, j);
    ip10 = I2D(ni, i+1, j);
    i0m1 = I2D(ni, i, j-1);
    i0p1 = I2D(ni, i, j+1);

    // evaluate derivatives
    d2tdx2 = temp_in[im10]-2*temp_in[i00]+temp_in[ip10];
    d2tdy2 = temp_in[i0m1]-2*temp_in[i00]+temp_in[i0p1];

    // update temperatures
    temp_out[i00] = temp_in[i00]+fact*(d2tdx2 + d2tdy2);
  }
}
\end{lstlisting}

\subsection{Compilation}
CUDA compilation was executed with the command:

\begin{verbatim}
    nvcc -O3 -arch=sm_50 cuda.cu -o cuda.exe -DDEBUG=0
\end{verbatim}

\textbf{Significato dei flag:}
\begin{itemize}
    \item \texttt{-O3}: maximum optimization.
    \item \texttt{-arch=sm\_50}: target GPU architecture.
    \item \texttt{-DDEBUG=0}: disables the DEBUG directives.
\end{itemize}

\subsection{Execution}
For the CUDA kernel, the chosen values for \texttt{threadsPerBlock} and \texttt{numBlocks} were the following.

\begin{lstlisting}
    dim3 threadsPerBlock(16, 16);
    dim3 numBlocks((ni + 15) / 16, (nj + 15) / 16);

    // ---- CUDA events start ----
    cudaEvent_t start, stop;
    cudaEventCreate(&start);
    cudaEventCreate(&stop);

    // Execute the modified version using same data
    printf("Execute GPU kernel...\n");
    cudaEventRecord(start, 0); // start timing

    for (istep = 0; istep < nstep; istep++) {
        step_kernel_mod <<<numBlocks, threadsPerBlock>>> (ni, nj, tfac, d_temp1, d_temp2);
        cudaDeviceSynchronize();

        cudaError_t err = cudaGetLastError();
        if (err != cudaSuccess)
            printf("CUDA Error: %s\n", cudaGetErrorString(err));
        
        float* d_tmp = d_temp1;
        d_temp1 = d_temp2;
        d_temp2 = d_tmp;
    }

    cudaEventRecord(stop, 0); // stop timing
    cudaEventSynchronize(stop);

\end{lstlisting}

\noindent
Each 16×16 thread block computes a portion of the grid, while the blocks are arranged to cover the entire matrix.

\noindent
Timing was performed using \texttt{cudaEvent\_t} to measure the execution time of the GPU kernel.

\subsection{Profiling}
Command: \texttt{nvprof build/cuda}

\hrule

\subsubsection*{950M:}
{\scriptsize
\begin{verbatim}
Execute the CPU-only reference version...
- Elapsed time: 0.188 s
Execute GPU kernel...
- Elapsed time: 140.515 ms
Max Error = 0.000008
The Max Error of 0.00001 is within ACCEPTABLE bounds.
- Total elapsed time: 2.460 s
==768== Profiling application: build/cuda
==768== Profiling result:
Type  Time(%)      Time     Calls       Avg       Min       Max  Name
GPU:   72.95%  202.16ms         2  101.08ms  2.7665ms  199.40ms  [CUDA memcpy HtoD]
       24.93%  69.101ms       200  345.50us  342.59us  355.94us  step_kernel_mod(int, int, float, float*, float*)
        2.12%  5.8774ms         2  2.9387ms  2.5962ms  3.2812ms  [CUDA memcpy DtoH]
API:   90.99%  1.69655s         2  848.27ms  699.40us  1.69585s  cudaMalloc
        7.10%  132.39ms       200  661.97us  423.90us  7.6760ms  cudaDeviceSynchronize
        0.96%  17.932ms         4  4.4831ms  3.1292ms  7.7131ms  cudaMemcpy
        0.45%  8.4117ms         2  4.2059ms  2.4000us  8.4093ms  cudaEventCreate
        0.43%  8.0093ms       200  40.046us  9.9000us  1.0583ms  cudaLaunchKernel
        0.04%  796.50us         2  398.25us  302.30us  494.20us  cudaFree
        0.01%  213.70us       200  1.0680us     200ns  31.300us  cudaGetLastError
        0.00%  90.700us         1  90.700us  90.700us  90.700us  cudaEventSynchronize
        0.00%  81.700us       101     808ns     200ns  31.300us  cuDeviceGetAttribute
        0.00%  30.100us         2  15.050us  10.700us  19.400us  cudaEventRecord
        0.00%  5.7000us         2  2.8500us  1.3000us  4.4000us  cudaEventDestroy
        0.00%  4.3000us         1  4.3000us  4.3000us  4.3000us  cuDeviceGetPCIBusId
        0.00%  3.9000us         1  3.9000us  3.9000us  3.9000us  cudaEventElapsedTime
        0.00%  2.3000us         3     766ns     400ns  1.4000us  cuDeviceGetCount
        0.00%  1.6000us         2     800ns     300ns  1.3000us  cuDeviceGet
        0.00%  1.5000us         1  1.5000us  1.5000us  1.5000us  cuDeviceGetName
        0.00%     800ns         1     800ns     800ns     800ns  cuDeviceTotalMem
        0.00%     400ns         1     400ns     400ns     400ns  cuDeviceGetUuid
\end{verbatim}
}

\hrule

\subsubsection*{T4:}
{\scriptsize
\begin{verbatim}
Execute the CPU-only reference version...
- Elapsed time: 36.761 s
Execute GPU kernel...
- Elapsed time: 1741.090 ms
Max Error = 0.000015
The Max Error of 0.00002 is within ACCEPTABLE bounds.
- Total elapsed time: 46.074 s
==5152== Profiling application: /content/drive/MyDrive/Colab Notebooks/HPC-CUDA/build/cuda
==5152== Profiling result:
Type  Time(%)      Time     Calls       Avg       Min       Max  Name
GPU:   67.30%  1.73103s       200  8.6551ms  8.0181ms  11.542ms  step_kernel_mod(int, int, float, float*, float*)
       16.79%  431.89ms         2  215.95ms  203.45ms  228.45ms  [CUDA memcpy HtoD]
       15.91%  409.13ms         2  204.56ms  201.12ms  208.01ms  [CUDA memcpy DtoH]
API:   62.34%  1.73679s       200  8.6840ms  8.0240ms  11.548ms  cudaDeviceSynchronize
       30.23%  842.24ms         4  210.56ms  201.46ms  228.69ms  cudaMemcpy
        7.16%  199.38ms         2  99.688ms  144.08us  199.23ms  cudaMalloc
        0.13%  3.7350ms       200  18.674us  5.4010us  429.51us  cudaLaunchKernel
        0.13%  3.6890ms         2  1.8445ms  870.86us  2.8181ms  cudaFree
        0.01%  168.47us       114  1.4770us     185ns  65.476us  cuDeviceGetAttribute
        0.00%  67.106us       200     335ns     105ns     674ns  cudaGetLastError
        0.00%  33.409us         2  16.704us  10.078us  23.331us  cudaEventRecord
        0.00%  27.898us         2  13.949us  1.1280us  26.770us  cudaEventCreate
        0.00%  12.296us         1  12.296us  12.296us  12.296us  cuDeviceGetName
        0.00%  6.8260us         1  6.8260us  6.8260us  6.8260us  cuDeviceGetPCIBusId
        0.00%  4.4120us         1  4.4120us  4.4120us  4.4120us  cudaEventSynchronize
        0.00%  3.6190us         1  3.6190us  3.6190us  3.6190us  cudaEventElapsedTime
        0.00%  3.3080us         2  1.6540us  1.1110us  2.1970us  cudaEventDestroy
        0.00%  1.6790us         3     559ns     233ns  1.2100us  cuDeviceGetCount
        0.00%     716ns         2     358ns     179ns     537ns  cuDeviceGet
        0.00%     529ns         1     529ns     529ns     529ns  cuDeviceTotalMem
        0.00%     372ns         1     372ns     372ns     372ns  cuModuleGetLoadingMode
        0.00%     302ns         1     302ns     302ns     302ns  cuDeviceGetUuid
\end{verbatim}
}

\hrule

\newpage
\subsection{Results}
The GPU simulation results were compared with the CPU version, and the speedup was calculated as:
\[
\text{speedup} = \frac{T_\text{CPU}}{T_\text{GPU}}
\]

\begin{table}[h!]
    \centering
    \begin{tabular}{lccc}
    \hline
    \textbf{GPU} & $\mathbf{T_{CPU}}$ (ms) & $\mathbf{T_{GPU}}$ (ms) & \textbf{Speedup} \\
    \hline
    950M & 34967 & 15086 & x2.32 \\
    T4   & 37178 & 1645  & x22.60 \\
    \hline
    \end{tabular}
\end{table}